\subsection{The matrix constrained-KP correspondence}
\label{sec:BA-constrained-KP}

The weighted BA identities also have a constrained-KP operator
interpretation for arbitrary positive $L,M,N$. We give the operator
dictionary and its higher-time normalization, then relate those times
to the canonical flows. The one-point specialization in
\cref{sec:BA-one-point-specialization} recovers scalar operators.

\subsubsection{Matrix operators from the BA bases}

For this operator comparison put $x=-it$ and
$\partial=\partial_x=i\partial_t$. Assemble the
ordinary and dual BA bases into columns
\[
 \boldsymbol\psi(P)=(\psi_1(P),\ldots,\psi_N(P))^{\trans},\qquad
 \boldsymbol\psi^\vee(P)=(\psi_1^\vee(P),\ldots,\psi_N^\vee(P))^{\trans}.
\]
The kernel rate in \cref{eq:general-kernel-rate} specifies the paired
primitive
\begin{equation}\label{eq:KP-paired-primitive}
 \partial^{-1}\!\left(
 (\boldsymbol\psi^\vee(Q))^{\trans}\boldsymbol\psi(P)\right)
 =\Omega(P,Q),
\end{equation}
with its integration constant fixed by the kernel normalization.
Define the diagonal puncture-constant matrices
$\mathsf C_A=\diag(a_{A,1},\ldots,a_{A,N})$, $A\in\{L,M\}$,
using \cref{eq:general-h-normalization}. The weighted equations give
\begin{equation}\label{eq:BA-constrained-operators}
\begin{aligned}
 \mathcal A_L&=\Id_N\partial+\mathsf C_L+
  2\sum_{\ell=1}^{L}\kappa_\ell\boldsymbol\psi(Q_\ell)
     \partial^{-1}(\boldsymbol\psi^\vee(Q_\ell))^{\trans},\\
 \mathcal A_M&=\Id_N\partial+\mathsf C_M+
  2\sum_{j=L+1}^{K}\kappa_j\boldsymbol\psi(Q_j)
     \partial^{-1}(\boldsymbol\psi^\vee(Q_j))^{\trans},
 \qquad K=L+M.
\end{aligned}
\end{equation}
The matrix size $N$ counts punctures; the two sums contain $L$ and
$M$ source columns, respectively. Each integral term pairs an input
column with the dual row, takes the
scalar primitive and multiplies by the ordinary column. Thus
\begin{equation}\label{eq:KP-matrix-spectral}
 \mathcal A_A\boldsymbol\psi(P)=h_A(P)\boldsymbol\psi(P),\qquad
 \mathcal A_A^{\rm ad}\boldsymbol\psi^\vee(P)
        =h_A(P)\boldsymbol\psi^\vee(P),\qquad A\in\{L,M\}.
\end{equation}
Here $\mathrm{ad}$ is formal transpose-adjoint, with
$\partial^{\rm ad}=-\partial$ and no complex conjugation.
The normalized expansions at all $N$ punctures determine an
invertible matrix wave operator. Both $\mathcal A_L,\mathcal A_M$
dress time-independent diagonal Laurent operators, so
\begin{equation}\label{eq:KP-matrix-commutation}
 [\mathcal A_L,\mathcal A_M]=0.
\end{equation}
Appendix~\ref{app:KP-normalization} supplies this formal argument and
the higher-time details.

Finite-rank integral operators of this form belong to established
matrix constrained-KP theory
\cite[equations (7)--(14)]{ChvartatskyiSydorenko2013}.
The scalar vector-constrained case, its rational reductions and
Grassmannian description are treated in
\cite[Theorems~4.1 and~5.1]{HelminckVanDeLeur1998} and
\cite{KricheverRational1995}. Residue-based BA self-consistency is
developed in \cite[Section~1]{KricheverSelfConsistent1995}.
The explicit dictionary here uses the two separate interpolation
groups and their puncture constants.

\subsubsection{Normalized KP times and source evolution}

The puncture constants determine how this operator form is normalized.
Define the diagonal gauge matrix $\mathsf G(x)=e^{-x\mathsf C_L}$ and put
$\widetilde{\boldsymbol\psi}=\mathsf G^{-1}\boldsymbol\psi$,
$\widetilde{\boldsymbol\psi}^{\vee}=\mathsf G^{\trans}\boldsymbol\psi^\vee$.
Then the first operator becomes
\begin{equation}\label{eq:KP-normalized-L}
 \mathcal L_{\rm KP}=\mathsf G^{-1}\mathcal A_L\mathsf G
 =\Id_N\partial+
 2\sum_{\ell=1}^{L}\kappa_\ell\widetilde{\boldsymbol\psi}(Q_\ell)
 \partial^{-1}(\widetilde{\boldsymbol\psi}^{\vee}(Q_\ell))^{\trans}.
\end{equation}
The second, $\mathcal M_{\rm KP}=\mathsf G^{-1}\mathcal A_M\mathsf G$,
retains the constant term
\begin{equation}\label{eq:KP-residual-phase}
 \mathsf C_M-\mathsf C_L
 =2\diag(b_1^N,\ldots,b_N^N).
\end{equation}
Thus the residual matrix is precisely the third-group phase data of
the canonical Hamiltonian. A common gauge removes both constant
terms, after additive shifts of $h_L,h_M$, only when their difference
is scalar; this holds automatically for $N=1$.

Use $h_L(P)$ as the local spectral coordinate at every puncture.
The normalized KP times $x_n$, with $x_1=x$, prescribe leading
ordinary factors $\exp(\sum_{n=1}^{n_*}x_n h_L(P)^n)e_k$ at $R_k$;
the dual factors have the opposite exponent. Here $n_*$ is any finite
number of active times and $e_k$ is the $k$th coordinate vector.
In the original coordinates,
\[
 \operatorname{pp}_{R_k}h_L^2=s_k^{-2}+2a_{L,k}s_k^{-1}.
\]
The regular constant $\operatorname{reg}_{R_k}h_L^2$ also supplies a
component phase. These terms explain why the second KP time
need not coincide with an uncorrected double-pole BA time.

\begin{proposition}[Constrained-KP evolution on the BA realization]
\label{prop:KP-hierarchy}
Let $\mathcal B_n=(\mathcal L_{\rm KP}^{\,n})_+$, where $+$ retains
the nonnegative differential powers. On a regular BA time domain,
\begin{equation}\label{eq:KP-source-flows}
\begin{aligned}
 \partial_{x_n}\widetilde{\boldsymbol\psi}(P)
   &=\mathcal B_n\widetilde{\boldsymbol\psi}(P),\\
 \partial_{x_n}\widetilde{\boldsymbol\psi}^{\vee}(P)
   &=-\mathcal B_n^{\rm ad}\widetilde{\boldsymbol\psi}^{\vee}(P),
 \qquad n\geq1.
\end{aligned}
\end{equation}
Both operators retain their finite-rank representations and satisfy
\begin{equation}\label{eq:KP-operator-flows}
 \partial_{x_n}\mathcal L_{\rm KP}
       =[\mathcal B_n,\mathcal L_{\rm KP}],\qquad
 \partial_{x_n}\mathcal M_{\rm KP}
       =[\mathcal B_n,\mathcal M_{\rm KP}].
\end{equation}
In particular, evaluation of \cref{eq:KP-source-flows} at the fixed
$Q_a$ gives the ordinary and dual source evolution for both groups.
\end{proposition}
\begin{proof}
In the $h_L$ local coordinates, the normalized matrix wave operator
dresses $\partial$ to $\mathcal L_{\rm KP}$. Subtracting the indicated
differential action from each BA time derivative leaves a section
with zero leading values at every puncture; BA uniqueness gives
\cref{eq:KP-source-flows}. The finite-rank commutator identity then
gives \cref{eq:KP-operator-flows}. The normalization and adjoint
arguments are detailed in Appendix~\ref{app:KP-normalization}.
\end{proof}

\subsubsection{Relation to the canonical mode flows}

The sources in \cref{eq:BA-constrained-operators} are precisely the
ordinary and dual evaluations used for the $v,w$ modes in
\cref{eq:multipoint-mode-arrays}. Evaluating the $L$ spectral problem
at the $M$-group points and the $M$ problem at the $L$-group points
recovers their weighted equations; the kernel derivative supplies
the $u$ equations. Multiwave equations involving eigenfunctions,
adjoint eigenfunctions and pair potentials also occur in
\cite[equations (57)--(58)]{AratynGomesNissimovPacheva2001}.

On each birational BA level, \cref{thm:BA-higher-flows} supplies
spectral Hamiltonians spanning all Abel directions. Matching a KP
velocity is therefore a finite linear problem. The remaining endpoint
phases are generated by the power charges in
\cref{eq:canonical-charges}. Each normalized KP time consequently
has a canonical spectral-plus-gauge generator on that level, with
its coefficients held fixed after choosing the level.
Appendix~\ref{app:KP-normalization} gives the construction, and
\cref{sec:numeric-KP-time} checks the second time in genus two.

This identifies the operators and their BA evolution. The pointwise
source arrays contain $NL+NM$ canonical pairs, whereas the full mode
system also retains $LM$ pair-composite pairs through its kernel
primitives. Identifying a KP field Poisson bracket with the complete
finite canonical bracket requires a further constraint calculation.
